\chapter{Objetivos}

\section{Objetivo geral}

Desenvolver e analisar uma sequência didática apoiada por uma plataforma interativa para o ensino introdutório de fundamentos de \textit{machine learning}, como conteúdo de Computação, a estudantes do Ensino Médio.

\section{Objetivos específicos}

Os objetivos específicos desta pesquisa compreendem:

\begin{itemize}
    \item analisar referências sobre ensino de inteligência artificial e \textit{machine learning} na Educação Básica, visualização interativa, múltiplas representações e avaliação da compreensão conceitual;
    \item relacionar os objetivos de aprendizagem da sequência aos eixos e às habilidades pertinentes da Computação na Educação Básica, sem supor que a norma prescreva algoritmos específicos;
    \item organizar uma sequência didática que trabalhe regressão linear, regressão logística, \textit{k}-vizinhos mais próximos, árvore de decisão e \textit{k}-means por meio do ciclo prever, experimentar, observar e explicar;
    \item ajustar cenários sintéticos, atividades de organização de tabelas pela turma, controles e explicações matemáticas da plataforma às condições do campo de aplicação;
    \item aplicar a sequência em uma turma de Ensino Médio em que o pesquisador é o professor;
    \item analisar, nas explicações, interpretações de gráficos e tabelas e escolhas justificadas dos estudantes, indícios de compreensão sobre origem dos dados, finalidade da tarefa, regra, parâmetro e resultado, presentes no diagnóstico inicial, no desafio final individual e nos registros docentes breves;
    \item sistematizar a plataforma, a sequência e os instrumentos de aplicação como produto educacional.
\end{itemize}
